% ---- pg-383 （书页 371） ----

(e) 用第一主成分作 Q-Q 图. 解释该图.

\noindent 8.23\quad 在阿拉斯加鱼类及游戏部工作的一位科学家研究灰熊, 旨在维持健康种群. 对 $n=61$ 头灰熊的测量数据综合出以下统计量:

\begin{mvtable}
\begin{tabular}{lcccccc}
\toprule
变量 & 重量(kg) & 体长(cm) & 颈围(cm) & 体围(cm) & 头长(cm) & 头宽(cm) \\
\midrule
样本均值 & 95.52 & 164.38 & 55.69 & 93.39 & 17.98 & 31.13 \\
\bottomrule
\end{tabular}
\end{mvtable}

\noindent 协方差矩阵
\[
\mathbf{S}=
\begin{bmatrix}
3\ 266.46 & 1\ 343.97 & 731.54 & 1\ 175.50 & 162.68 & 238.37\\
1\ 343.97 & 721.91 & 324.25 & 537.35 & 80.17 & 117.73\\
731.54 & 324.25 & 179.28 & 281.17 & 39.15 & 56.80\\
1\ 175.50 & 537.35 & 281.17 & 474.98 & 63.73 & 94.85\\
162.68 & 80.17 & 39.15 & 63.73 & 9.95 & 13.88\\
238.37 & 117.73 & 56.80 & 94.85 & 13.88 & 21.26
\end{bmatrix}
\]

(a) 用协方差矩阵作主成分分析. 考察这组数据能否用低于 6 维的结构来有效地描述?

(b) 用相关矩阵作主成分分析.

(c) 评价两种分析结果的相似与相异之处.

\noindent 8.24\quad 对\hyperref[tab:5.8]{表 5.8} 中的数据, 添加变量 $x_6=$ 正常加班时间, 其数据如下(按横排顺序):
\[
\begin{array}{rrrrrrrr}
6\ 187 & 7\ 336 & 6\ 988 & 6\ 964 & 8\ 425 & 6\ 778 & 5\ 922 & 7\ 307\\
7\ 679 & 8\ 259 & 10\ 954 & 9\ 353 & 6\ 291 & 4\ 969 & 4\ 825 & 6\ 019
\end{array}
\]
对添加了变量的数据按例 8.10 的做法重作一遍.

\noindent 8.25\quad 用例 8.10 中的警察超时工作时间数据, 基于 $d_{Uj}^2$ 的和构造另一种控制图, 用以监控在原始数据中其他主成分所综合的变差.

\noindent 8.26\quad 考虑\hyperref[tab:4.6]{表 4.6} 中的心理刻画数据. 用五个变量(独立性、支持精神、仁爱心、坚定性和领导力)分布对协方差阵 $\mathbf{S}$ 和相关阵 $\mathbf{R}$ 作主成分分析. 分析应包含下列内容:

(a) 确定适当数量的主成分以有效地综合变差. 构造碎石图来帮助你确定维数.

(b) 解释样本主成分.

(c) 用前两个主成分的值, 作数据的散点图. 用 $\hat{y}_1$ 作纵轴, $\hat{y}_2$ 作横轴. 你能区分两个社会经济水平组和(或)两个性别组吗? 是否存在离群值?

(d) 对协方差矩阵的第一特征值 $\lambda_1$ 构造 95\% 的置信区间.

\noindent 8.27\quad \hyperref[tab:7.7]{表 7.7} 给出了纸浆与纸性的数据. 用纸性的四个变量(BL, EM, SF 和 BS)分别基于协方差阵 $\mathbf{S}$ 和相关阵 $\mathbf{R}$ 作主成分分析. 分析应包含下列内容:

(a) 确定适当数量的主成分以有效地综合变差. 构造碎石图来帮助你确定维数.

(b) 解释样本主成分.

(c) 你是否认为可能构造一个“纸强度”指标, 有效包含四个纸变量的信息? 作出解释.

(d) 用前两个主成分的值, 作数据的散点图. 用 $\hat{y}_1$ 作纵轴, $\hat{y}_2$ 作横轴. 识别数据中的离群值.

\noindent 8.28\quad 为评估马里(西非)斯卡索地区自然资源是否可支持扩大的食物需求的安全线, 一项研究得到一组调查数据. 总共调查了 $n=76$ 个农夫, 得到了关于 9 个变量的观测. 变量罗列如下:
